%! TEX root = main.tex
% Stile del documento
\documentclass[12pt,openany]{amsbook}
\usepackage[T1]{fontenc}
\usepackage[utf8]{inputenc}
\usepackage[a4paper]{geometry}
\geometry{verbose,tmargin=1.8cm,bmargin=2cm,lmargin=2cm,rmargin=2cm,headheight=1.8cm,headsep=1cm,footskip=1cm}
% Localizzazione (traduce nomi automatici nella lingua selezionata)
\usepackage[italian]{babel}
% Altri pacchetti per la formattazione
\usepackage{lmodern}
\usepackage{hyphenat}
\usepackage{amsmath}
\usepackage{amsthm}
\usepackage{amssymb}
\usepackage{mathrsfs}
\usepackage{mathtools}
\usepackage{cancel}
\usepackage{enumerate}
\usepackage{graphicx}
 % Impostazioni dei metadati del PDF
\usepackage[unicode=true, bookmarks=true, bookmarksnumbered=true, bookmarksopen=true, bookmarksopenlevel=1, breaklinks=false, pdfborder={0 0 0}, pdfborderstyle={}, backref=page, colorlinks=false]
{hyperref}
\hypersetup{
    pdftitle={Eserciziario di Sicurezza delle Reti},
    pdfauthor={Lorenzo Prosseda},
    pdfsubject={Corso di Sicurezza delle reti del prof. Bregni, Politecnico di Milano 2019\---2020},
    pdfkeywords={sicurezza, reti, crittografia}
}
% Percorso della cartella contenente le immagini
\graphicspath{ {../img/} }
% Impostazioni di stile per gli ambienti teorema, definizione, esempio, ecc
\theoremstyle{plain}
\newtheorem{thm}{Teorema}[section] % resetta numerazione ad ogni capitolo
\renewcommand{\thethm}{\arabic{chapter}.\arabic{section}.\arabic{thm}} % numerazione ambienti teorema
\numberwithin{equation}{section} % numerazione equazioni
% Definizione ambienti personalizzati
\theoremstyle{definition}
\newtheorem{defn}[thm]{Definizione}
\newtheorem{exmp}[thm]{Esempio}
\newtheorem{obsv}[thm]{Osservazione}
\theoremstyle{remark}
\newtheorem{note}[thm]{Nota}
% Numerazione personalizzata sezioni e sottosezioni
\renewcommand\thesection{\thechapter.\arabic{section}}
\renewcommand\thesubsection{\thesection.\arabic{subsection}}
% Comando trattino "non spezzante" evita sillabazione sul trattino
\newcommand\nbdash{\nobreakdash-\hspace{0pt}}
% Non spezzare e seguenti parole:
\hyphenation{Definizione}
\hyphenation{Osservazione}
\hyphenation{Equazione}
\hyphenation{Teorema}
\hyphenation{Nota}
% Solleva posizione lettera greca chi minuscola
\usepackage{letltxmacro}
\LetLtxMacro{\lowchi}{\chi}
\renewcommand\chi{\text{\raisebox{2pt}{$\lowchi$}}}

% Comando per ambiente problema e soluzione
\newcounter{problem}
\newcounter{solution}
\newcommand\Problem[1]{%
  \stepcounter{problem}%
  \bigskip
  \textbf{Esercizio \theproblem.}~\emph{(#1)}~%
  \setcounter{solution}{0}%
}
\newcommand\TheSolution{%
  \textbf{Soluzione}\\%
}
\newcommand\ASolution[1]{%
  \stepcounter{solution}%
  \bigskip
  \textbf{Soluzione \thesolution:}\hfill\emph{(#1)}\\%
}

% In questo ambiente risiede il contenuto del documento
\begin{document}
    % Pagina del titolo
    %! TEX root = main.tex

\begin{titlepage}
    \centering % centra tutto in questo ambiente
    \vspace*{\stretch{0.8}}
    \Huge\textbf{Esercizi e problemi di\\Sicurezza delle Reti}\\
    \vspace*{\stretch{0.6}}
    \Large\emph{Lorenzo Prosseda}\\
    \bigskip
    \large a.a. 2019\---2020 \\
    \vspace*{\stretch{5.0}}
    \begin{figure}[h] % inserisci logo ed estratto della licenza GFDL
        \includegraphics[width=2cm]{gfdl-logo.pdf}
        \bigskip \\
        \small\emph{Copyright \textcopyright\,2019 Lorenzo Prosseda. Permission is granted to copy, distribute and/or modify this document under the terms of the GNU Free Documentation License, Version 1.3 or any later version published by the Free Software Foundation; with no Invariant Sections, no Front-Cover Texts, and no Back-Cover Texts. A copy of the license is included in the file called ``LICENSE''.}
    \end{figure}
    \vspace*{\stretch{1.0}}
\end{titlepage}


    % Genera indice di capitoli e sezioni
    \tableofcontents

    % Includi file dei capitoli
    \part{Esercizi}
    %! TEX root = main.tex
% Capitolo 1

\chapter{Firma digitale}

    \bigskip
    \section{Firma RSA}
        \Problem{27/02/2013}
        Bob adotta il sistema di firma elettronica RSA; egli pubblica il modulo $n=437$ e l'esponente di 
        cifratura $e=17$.
        \begin{enumerate}
            \item Verificare la correttezza dei dati forniti in base alle ipotesi di del metodo RSA,
            \item Calcolare la firma $A$ del messaggio $P=77$.
            \item Il valore $A\equiv\,2\,(\mathrm{mod}\,n)$ è una firma valida per quale messaggio $P$?
        \end{enumerate}

    \bigskip
    \section{Firma Cieca}
        \Problem{25/09/2013}
        Alice e Bob adottano il sistema di firma cieca RSA; Alice pubblica il modulo $n=391$ e l'esponente 
        di cifrature $e=13$. Bob estrae il numero casuale segreto (nonce) $k=12$ e chiede ad Alice di firmare 
        alla cieca il messaggio $P=32$.
        \begin{enumerate}
            \item Verificare la correttezza dei dati forniti secondo le ipotesi del metodo di firma cieca RSA.
            \item Calcolare i messaggi scambiati da Alice e Bob, e la firma $A$ del messaggio $P$.
        \end{enumerate}

        \Problem{18/07/2014}
        Alice e Bob adottano il sistema di firma cieca RSA; Alice pubblica il modulo $n=319$ e l'esponente 
        di cifrature $e=17$. Bob estrae il numero casuale segreto (nonce) $k=25$ e chiede ad Alice di firmare 
        alla cieca il messaggio $P=299$.
        \begin{enumerate}
            \item Verificare la correttezza dei dati forniti secondo le ipotesi del metodo di firma cieca RSA.
            \item Calcolare i messaggi scambiati da Alice e Bob, e la firma $A$ del messaggio $P$.
        \end{enumerate}

        \Problem{03/07/2015}
        Alice e Bob adottano il sistema di firma cieca RSA; Alice pubblica il modulo $n=551$ e l'esponente 
        di cifrature $e=17$. Bob estrae il numero casuale segreto (nonce) $k=12$ e chiede ad Alice di firmare 
        alla cieca il messaggio $P=500$.
        \begin{enumerate}
            \item Verificare la correttezza dei dati forniti secondo le ipotesi del metodo di firma cieca RSA.
            \item Calcolare i messaggi scambiati da Alice e Bob, e la firma $A$ del messaggio $P$.
        \end{enumerate}

        \Problem{21/07/2016}
        Alice e Bob adottano il sistema di firma cieca RSA; Alice pubblica il modulo $n=899$ e l'esponente 
        di cifrature $e=17$. Bob estrae il numero casuale segreto (nonce) $k=501$ e chiede ad Alice di firmare 
        alla cieca il messaggio $P=500$.
        \begin{enumerate}
            \item Verificare la correttezza dei dati forniti secondo le ipotesi del metodo di firma cieca RSA.
            \item Calcolare i messaggi scambiati da Alice e Bob, e la firma $A$ del messaggio $P$.
        \end{enumerate}

        \Problem{09/02/2017}
        Alice e Bob adottano il sistema di firma cieca RSA; Alice pubblica il modulo $n=1073$ e l'esponente 
        di cifrature $e=25$. Bob estrae il numero casuale segreto (nonce) $k=17$ e chiede ad Alice di firmare 
        alla cieca il messaggio $P=50$.
        \begin{enumerate}
            \item Verificare la correttezza dei dati forniti secondo le ipotesi del metodo di firma cieca RSA.
            \item Calcolare i messaggi scambiati da Alice e Bob, e la firma $A$ del messaggio $P$.
        \end{enumerate}

        \Problem{30/08/2018}
        Alice e Bob adottano il sistema di firma cieca RSA; Alice pubblica il modulo $n=667$ e l'esponente 
        di cifrature $e=39$. Bob estrae il numero casuale segreto (nonce) $k=20$ e chiede ad Alice di firmare 
        alla cieca il messaggio $P=500$.
        \begin{enumerate}
            \item Verificare la correttezza dei dati forniti secondo le ipotesi del metodo di firma cieca RSA.
            \item Calcolare i messaggi scambiati da Alice e Bob, e la firma $A$ del messaggio $P$.
        \end{enumerate}

    \bigskip
    \section{Firma di El Gamal}
        \Problem{08/02/2013}
        Bob adotta il sistema di firma elettronica di El Gamal e pubblica le seguenti informazioni: \[
            p=113,\, \alpha = 7,\, \beta = \alpha^a\,\mathrm{mod}\,p
        ;\] inoltre Bob mantiene segreto l'esponente $a=29$.
        \begin{enumerate}
            \item Verificare la correttezza dei dati forniti, in base alle ipotesi del metodo di El Gamal. 
                Se $\alpha=7$ non risultasse una scelta valida, Bob userà invece $\alpha=6$ (da verificare); 
                se nemmeno questa scelta risultasse valida, Bob rinuncerà a proseguire (l'esercizio termina); 
                calcolare $\beta$.
            \item Bob estrae il numero casuale segreto (nonce) $k=51$; per questo valore di $k$, calcolare 
                la firma $A$ del messaggio $P=101$.
            \item Verificare se anche la firma $A^{\prime}=(r^{\prime}, s^{\prime})=(103,46)$ è valida per 
                lo stesso messaggio $P=101$.
        \end{enumerate}

        \Problem{11/07/2013}
        Bob adotta il sistema di firma elettronica di El Gamal e pubblica le seguenti informazioni: \[
            p=113,\, \alpha = 6,\, \beta = \alpha^a\,\mathrm{mod}\,p = 37
        ;\] inoltre Bob mantiene segreto l'esponente $a\in (1,\,p-2]$.
        \begin{enumerate}
            \item Verificare la correttezza dei dati forniti, in base alle ipotesi del metodo di El Gamal. 
                Se $\alpha=6$ non risultasse una scelta valida, Bob userà invece $\alpha=7$ (da verificare); 
                se nemmeno questa scelta risultasse valida, Bob rinuncerà a proseguire (l'esercizio termina).
            \item Bob estrae il numero casuale segreto (nonce) $k\perp p-1$; usando due volte lo stesso $k$, 
                Bob calcola le firme $A_1=(r_1,\,s_1)=(55,\,62)$ e $A_2=(r_2,\,s_2)=(55,\,39)$, rispetto ai 
                messaggi $P_1=7$ e $P_2=48$. Oscar intercetta $A_1$ e $A_2$; calcolare $k$ e $a$ sfruttando 
                l'attacco del nonce ripetuto.
        \end{enumerate}

        \Problem{13/09/2013}
        Bob adotta il sistema di firma elettronica di El Gamal e pubblica le seguenti informazioni: \[
            p=113,\, \alpha = 6,\, \beta = \alpha^a\,\mathrm{mod}\,p=31
        ;\] inoltre Bob mantiene segreto l'esponente $a\in (1,\,p-2]$.
        \begin{enumerate}
            \item Bob estrae il numero casuale segreto (nonce) $k\perp p-1$; usando lo stesso $k$, 
                Bob calcola le seguenti firme $A_k$ per i messaggi $P_k$:
                \[\begin{array}{ll}
                    A_1=(r_1,\,s_1)=(92,\,4) & P_1=12\\
                    A_2=(r_2,\,s_2)=(92,\,44) & P_2=100\\
                    A_3=(r_3,\,s_3)=(92,\,12) & P_3=35
                \end{array}\]
                Verificare la validità delle tre firme.
            \item Oscar intercetta i tre messaggi $(P_k,\,A_k)$; sulla base delle sole firme verificate valide 
                (nel precedente quesito), calcolare $k$ e $a$ sfruttando l'attacco del nonce ripetuto.
        \end{enumerate}

        \Problem{03/03/2014}
        Bob adotta il sistema di firma elettronica di El Gamal e pubblica le seguenti informazioni: \[
            p=109,\, \alpha = 6,\, \beta = \alpha^a\,\mathrm{mod}\,p=48
        ;\] inoltre Bob mantiene segreto l'esponente $a\in (1,\,p-2]$.
        \begin{enumerate}
            \item Bob estrae il numero casuale segreto (nonce) $k\perp p-1$; usando lo stesso $k$, 
                Bob calcola le seguenti firme $A_k$ per i messaggi $P_k$:
                \[\begin{array}{ll}
                    A_1=(r_1,\,s_1)=(42,\,74) & P_1=106\\
                    A_2=(r_2,\,s_2)=(42,\,99) & P_2=10\\
                    A_3=(r_3,\,s_3)=(42,\,28) & P_3=56
                \end{array}\]
                Verificare la validità delle tre firme.
            \item Oscar intercetta i tre messaggi $(P_k,\,A_k)$; sulla base delle sole firme verificate valide 
                (nel precedente quesito), calcolare $k$ e $a$ sfruttando l'attacco del nonce ripetuto.
        \end{enumerate}

        \Problem{04/07/2014}
        Bob adotta il sistema di firma elettronica di El Gamal e pubblica le seguenti informazioni: \[
            p=181,\, \alpha = 7,\, \beta = \alpha^a\,\mathrm{mod}\,p
        ;\] inoltre Bob mantiene segreto l'esponente $a=56$.
        \begin{enumerate}
            \item Verificare la correttezza dei dati forniti, in base alle ipotesi del metodo di El Gamal. 
                Se $\alpha=7$ non risultasse una scelta valida, Bob userà invece $\alpha=10$ (da verificare); 
                se nemmeno questa scelta risultasse valida, Bob rinuncerà a proseguire (l'esercizio termina); 
                calcolare $\beta$.
            \item Bob estrae il numero casuale segreto (nonce) $k=23$; per questo valore di $k$, calcolare 
                la firma $A$ del messaggio $P=11$.
            \item Verificare se anche la firma $A^{\prime}=(r^{\prime}, s^{\prime})=(112,57)$ è valida per 
                lo stesso messaggio $P=11$.
        \end{enumerate}

        \Problem{17/09/2014}
        Bob adotta il sistema di firma elettronica di El Gamal e pubblica le seguenti informazioni: \[
            p=149,\, \alpha = 8,\, \beta = \alpha^a\,\mathrm{mod}\,p=127
        ;\] inoltre Bob mantiene segreto l'esponente $a\in (1,\,p-2]$.
        \begin{enumerate}
            \item Bob estrae il numero casuale segreto (nonce) $k\perp p-1$; usando lo stesso $k$, 
                Bob calcola le seguenti firme $A_k$ per i messaggi $P_k$:
                \[\begin{array}{ll}
                    A_1=(r_1,\,s_1)=(75,\,2) & P_1=10\\
                    A_2=(r_2,\,s_2)=(75,\,124) & P_2=20\\
                    A_3=(r_3,\,s_3)=(75,\,94) & P_3=30
                \end{array}\]
                Verificare la validità delle tre firme.
            \item Oscar intercetta i tre messaggi $(P_k,\,A_k)$; sulla base delle sole firme verificate valide 
                (nel precedente quesito), calcolare $k$ e $a$ sfruttando l'attacco del nonce ripetuto.
        \end{enumerate}

        \Problem{13/02/2015}
        Bob adotta il sistema di firma elettronica di El Gamal e pubblica le seguenti informazioni: \[
            p=107,\, \alpha = 5,\, \beta = \alpha^a\,\mathrm{mod}\,p=66
        ;\] inoltre Bob mantiene segreto l'esponente $a\in (1,\,p-2]$.
        \begin{enumerate}
            \item Bob estrae il numero casuale segreto (nonce) $k\perp p-1$; usando lo stesso $k$, 
                Bob calcola le seguenti firme $A_k$ per i messaggi $P_k$:
                \[\begin{array}{ll}
                    A_1=(r_1,\,s_1)=(84,\,74) & P_1=14\\
                    A_2=(r_2,\,s_2)=(84,\,23) & P_2=15\\
                    A_3=(r_3,\,s_3)=(84,\,40) & P_3=16
                \end{array}\]
                Verificare la validità delle tre firme.
            \item Oscar intercetta i tre messaggi $(P_k,\,A_k)$; sulla base delle sole firme verificate valide 
                (nel precedente quesito), calcolare $k$ e $a$ sfruttando l'attacco del nonce ripetuto.
        \end{enumerate}

        \Problem{02/03/2015}
        Bob adotta il sistema di firma elettronica di El Gamal e pubblica le seguenti informazioni: \[
            p=199,\, \alpha = 7,\, \beta = \alpha^a\,\mathrm{mod}\,p
        ;\] inoltre Bob mantiene segreto l'esponente $a=56$.
        \begin{enumerate}
            \item Verificare la correttezza dei dati forniti, in base alle ipotesi del metodo di El Gamal. 
                Se $\alpha=7$ non risultasse una scelta valida, Bob userà invece $\alpha=6$ (da verificare); 
                se nemmeno questa scelta risultasse valida, Bob rinuncerà a proseguire (l'esercizio termina); 
                calcolare $\beta$.
            \item Bob estrae il numero casuale segreto (nonce) $k=151$; per questo valore di $k$, calcolare 
                la firma $A$ del messaggio $P=200$.
            \item Verificare se anche la firma $A^{\prime}=(r^{\prime}, s^{\prime})=(44,100)$ è valida per 
                lo stesso messaggio $P=200$.
        \end{enumerate}

        \Problem{17/07/2015}
        Bob adotta il sistema di firma elettronica di El Gamal e pubblica le seguenti informazioni: \[
            p=107,\, \alpha = 7,\, \beta = \alpha^a\,\mathrm{mod}\,p=72
        ;\] inoltre Bob mantiene segreto l'esponente $a\in (1,\,p-2]$.
        \begin{enumerate}
            \item Bob estrae il numero casuale segreto (nonce) $k\perp p-1$; usando lo stesso $k$, 
                Bob calcola le seguenti firme $A_k$ per i messaggi $P_k$:
                \[\begin{array}{ll}
                    A_1=(r_1,\,s_1)=(67,\,58) & P_1=61\\
                    A_2=(r_2,\,s_2)=(67,\,31) & P_2=62\\
                    A_3=(r_3,\,s_3)=(67,\,100) & P_3=63
                \end{array}\]
                Verificare la validità delle tre firme.
            \item Oscar intercetta i tre messaggi $(P_k,\,A_k)$; sulla base delle sole firme verificate valide 
                (nel precedente quesito), calcolare $k$ e $a$ sfruttando l'attacco del nonce ripetuto.
        \end{enumerate}

        \Problem{04/07/2016}
        Bob adotta il sistema di firma elettronica di El Gamal e pubblica le seguenti informazioni: \[
            p=103,\, \alpha = 5,\, \beta = \alpha^a\,\mathrm{mod}\,p=84
        ;\] inoltre Bob mantiene segreto l'esponente $a\in (1,\,p-2]$.
        \begin{enumerate}
            \item Bob estrae il numero casuale segreto (nonce) $k\perp p-1$; usando lo stesso $k$, 
                Bob calcola le seguenti firme $A_k$ per i messaggi $P_k$:
                \[\begin{array}{ll}
                    A_1=(r_1,\,s_1)=(45,\,59) & P_1=50\\
                    A_2=(r_2,\,s_2)=(45,\,91) & P_2=51\\
                    A_3=(r_3,\,s_3)=(45,\,19) & P_3=52
                \end{array}\]
                Verificare la validità delle tre firme.
            \item Oscar intercetta i tre messaggi $(P_k,\,A_k)$; sulla base delle sole firme verificate valide 
                (nel precedente quesito), calcolare $k$ e $a$ sfruttando l'attacco del nonce ripetuto.
        \end{enumerate}

        \Problem{21/07/2016}
        Bob adotta il sistema di firma elettronica di El Gamal e pubblica le seguenti informazioni: \[
            p=109,\, \alpha = 10,\, \beta = \alpha^a\,\mathrm{mod}\,p=59
        ;\] inoltre Bob mantiene segreto l'esponente $a\in (1,\,p-2]$.
        \begin{enumerate}
            \item Bob estrae il numero casuale segreto (nonce) $k\perp p-1$; usando lo stesso $k$, 
                Bob calcola le seguenti firme $A_k$ per i messaggi $P_k$:
                \[\begin{array}{ll}
                    A_1=(r_1,\,s_1)=(79,\,10) & P_1=32\\
                    A_2=(r_2,\,s_2)=(79,\,20) & P_2=33\\
                    A_3=(r_3,\,s_3)=(79,\,102) & P_3=35
                \end{array}\]
                Verificare la validità delle tre firme.
            \item Oscar intercetta i tre messaggi $(P_k,\,A_k)$; sulla base delle sole firme verificate valide 
                (nel precedente quesito), calcolare $k$ e $a$ sfruttando l'attacco del nonce ripetuto.
        \end{enumerate}

        \Problem{09/09/2016}
        Bob adotta il sistema di firma elettronica di El Gamal e pubblica le seguenti informazioni: \[
            p=211,\, \alpha = 6,\, \beta = \alpha^a\,\mathrm{mod}\,p
        ;\] inoltre Bob mantiene segreto l'esponente $a=99$.
        \begin{enumerate}
            \item Verificare la correttezza dei dati forniti, in base alle ipotesi del metodo di El Gamal. 
                Se $\alpha=6$ non risultasse una scelta valida, Bob userà invece $\alpha=7$ (da verificare); 
                se nemmeno questa scelta risultasse valida, Bob rinuncerà a proseguire (l'esercizio termina); 
                calcolare $\beta$.
            \item Bob estrae il numero casuale segreto (nonce) $k=191$; per questo valore di $k$, calcolare 
                la firma $A$ del messaggio $P=200$.
            \item Verificare se anche la firma $A^{\prime}=(r^{\prime}, s^{\prime})=(167,157)$ è valida per 
                lo stesso messaggio $P=200$; se è valida, calcolare il valore di $k$ per cui è stata ricavata da Bob.
        \end{enumerate}

        \Problem{06/07/2018}
        Bob adotta il sistema di firma elettronica di El Gamal e pubblica le seguenti informazioni: \[
            p=127,\, \alpha = 2,\, \beta = \alpha^a\,\mathrm{mod}\,p
        ;\] inoltre Bob mantiene segreto l'esponente $a=65$.
        \begin{enumerate}
            \item Verificare la correttezza dei dati forniti, in base alle ipotesi del metodo di El Gamal. 
                Se $\alpha=2$ non risultasse una scelta valida, Bob userà invece $\alpha=\in\{3,\,4,\,5\}$ 
                (da verificare); se nemmeno queste scelte risultassero valide, Bob rinuncerà a proseguire 
                (l'esercizio termina); calcolare $\beta$.
            \item Bob estrae il numero casuale segreto (nonce) $k=19$; per questo valore di $k$, calcolare 
                la firma $A$ del messaggio $P=19$.
            \item Verificare se anche la firma $A^{\prime}=(r^{\prime}, s^{\prime})=(6,7)$ è valida per 
                lo stesso messaggio $P=19$; se è valida, calcolare il valore di $k$ per cui è stata ricavata da Bob.
        \end{enumerate}

        \Problem{27/07/2018}
        Bob adotta il sistema di firma elettronica di El Gamal e pubblica le seguenti informazioni: \[
            p=157,\, \alpha = 5,\, \beta = \alpha^a\,\mathrm{mod}\,p=109
        ;\] inoltre Bob mantiene segreto l'esponente $a\in (1,\,p-2]$.
        \begin{enumerate}
            \item Bob estrae il numero casuale segreto (nonce) $k\perp p-1$; usando lo stesso $k$, 
                Bob calcola le seguenti firme $A_k$ per i messaggi $P_k$:
                \[\begin{array}{ll}
                    A_1=(r_1,\,s_1)=(26,\,34) & P_1=10\\
                    A_2=(r_2,\,s_2)=(26,\,2) & P_2=11\\
                    A_3=(r_3,\,s_3)=(26,\,20) & P_3=12
                \end{array}\]
                Verificare la validità delle tre firme.
            \item Oscar intercetta i tre messaggi $(P_k,\,A_k)$; sulla base delle sole firme verificate valide 
                (nel precedente quesito), calcolare $k$ e $a$ sfruttando l'attacco del nonce ripetuto.
        \end{enumerate}

        \Problem{23/01/2019}
        Bob adotta il sistema di firma elettronica di El Gamal e pubblica le seguenti informazioni: \[
            p=113,\, \alpha = 6,\, \beta = \alpha^a\,\mathrm{mod}\,p=16
        ;\] inoltre Bob mantiene segreto l'esponente $a\in (1,\,p-2]$.
        \begin{enumerate}
            \item Bob estrae il numero casuale segreto (nonce) $k\perp p-1$; usando lo stesso $k$, 
                Bob calcola le seguenti firme $A_k$ per i messaggi $P_k$:
                \[\begin{array}{ll}
                    A_1=(r_1,\,s_1)=(94,\,13) & P_1=11\\
                    A_2=(r_2,\,s_2)=(94,\,10) & P_2=13\\
                    A_3=(r_3,\,s_3)=(94,\,9) & P_3=15
                \end{array}\]
                Verificare la validità delle tre firme.
            \item Oscar intercetta i tre messaggi $(P_k,\,A_k)$; sulla base delle sole firme verificate valide 
                (nel precedente quesito), calcolare $k$ e $a$ sfruttando l'attacco del nonce ripetuto.
        \end{enumerate}

        \Problem{23/01/2019}
        Bob adotta il sistema di firma elettronica di El Gamal e pubblica le seguenti informazioni: \[
            p=137,\, \alpha = 2,\, \beta = \alpha^a\,\mathrm{mod}\,p
        ;\] inoltre Bob mantiene segreto l'esponente $a=2$.
        \begin{enumerate}
            \item Verificare la correttezza dei dati forniti, in base alle ipotesi del metodo di El Gamal. 
                Se $\alpha=2$ non risultasse una scelta valida, Bob userà invece $\alpha\in \{3,\,4\}$ 
                (da verificare); se nemmeno questa scelta risultasse valida, Bob rinuncerà a proseguire 
                (l'esercizio termina); calcolare $\beta$.
            \item Bob estrae il numero casuale segreto (nonce) $k=21$; per questo valore di $k$, calcolare 
                la firma $A$ del messaggio $P=100$.
            \item Verificare se anche la firma $A^{\prime}=(r^{\prime}, s^{\prime})=(48,4)$ è valida per 
                lo stesso messaggio $P=100$; se è valida, calcolare il valore di $k$ per cui è stata ricavata da Bob.
        \end{enumerate}

        \Problem{11/02/2019}
        Bob adotta il sistema di firma elettronica di El Gamal e pubblica le seguenti informazioni: \[
            p=157,\, \alpha = 6,\, \beta = \alpha^a\,\mathrm{mod}\,p=37
        ;\] inoltre Bob mantiene segreto l'esponente $a\in (1,\,p-2]$.
        \begin{enumerate}
            \item Bob estrae il numero casuale segreto (nonce) $k\perp p-1$; usando lo stesso $k$, 
                Bob calcola le seguenti firme $A_k$ per i messaggi $P_k$:
                \[\begin{array}{ll}
                    A_1=(r_1,\,s_1)=(26,\,32) & P_1=20\\
                    A_2=(r_2,\,s_2)=(26,\,28) & P_2=24\\
                    A_3=(r_3,\,s_3)=(26,\,28) & P_3=28
                \end{array}\]
                Verificare la validità delle tre firme.
            \item Oscar intercetta i tre messaggi $(P_k,\,A_k)$; sulla base delle sole firme verificate valide 
                (nel precedente quesito), calcolare $k$ e $a$ sfruttando l'attacco del nonce ripetuto.
        \end{enumerate}

        \Problem{02/07/2019}
        Bob adotta il sistema di firma elettronica di El Gamal e pubblica le seguenti informazioni: \[
            p=137,\, \alpha = 4,\, \beta = \alpha^a\,\mathrm{mod}\,p
        ;\] inoltre Bob mantiene segreto l'esponente $a=129$.
        \begin{enumerate}
            \item Verificare la correttezza dei dati forniti, in base alle ipotesi del metodo di El Gamal. 
                Se $\alpha=4$ non risultasse una scelta valida, Bob userà invece $\alpha\in \{5,\,7\}$ 
                (da verificare); se nemmeno questa scelta risultasse valida, Bob rinuncerà a proseguire 
                (l'esercizio termina); calcolare $\beta$.
            \item Bob estrae il numero casuale segreto (nonce) $k=29$; per questo valore di $k$, calcolare 
                la firma $A$ del messaggio $P=101$.
            \item Verificare se anche la firma $A^{\prime}=(r^{\prime}, s^{\prime})=(117,32)$ è valida per 
                lo stesso messaggio $P=101$; se è valida, calcolare il valore di $k$ per cui è stata ricavata da Bob.
        \end{enumerate}

        \Problem{26/07/2019}
        Bob adotta il sistema di firma elettronica di El Gamal e pubblica le seguenti informazioni: \[
            p=131,\, \alpha = 6,\, \beta = \alpha^a\,\mathrm{mod}\,p=2
        ;\] inoltre Bob mantiene segreto l'esponente $a\in (1,\,p-2]$.
        \begin{enumerate}
            \item Bob estrae il numero casuale segreto (nonce) $k\perp p-1$; usando lo stesso $k$, 
                Bob calcola le seguenti firme $A_k$ per i messaggi $P_k$:
                \[\begin{array}{ll}
                    A_1=(r_1,\,s_1)=(30,\,90) & P_1=10\\
                    A_2=(r_2,\,s_2)=(30,\,45) & P_2=15\\
                    A_3=(r_3,\,s_3)=(30,\,1) & P_3=20
                \end{array}\]
                Verificare la validità delle tre firme.
            \item Oscar intercetta i tre messaggi $(P_k,\,A_k)$; sulla base delle sole firme verificate valide 
                (nel precedente quesito), calcolare $k$ e $a$ sfruttando l'attacco del nonce ripetuto.
        \end{enumerate}

        \Problem{31/08/2019}
        Bob adotta il sistema di firma elettronica di El Gamal e pubblica le seguenti informazioni: \[
            p=127,\, \alpha = 4,\, \beta = \alpha^a\,\mathrm{mod}\,p
        ;\] inoltre Bob mantiene segreto l'esponente $a=129$.
        \begin{enumerate}
            \item Verificare la correttezza dei dati forniti, in base alle ipotesi del metodo di El Gamal. 
                Se $\alpha=4$ non risultasse una scelta valida, Bob userà invece $\alpha\in \{5,\,6\}$ 
                (da verificare); se nemmeno questa scelta risultasse valida, Bob rinuncerà a proseguire 
                (l'esercizio termina); calcolare $\beta$.
            \item Bob estrae il numero casuale segreto (nonce) $k=25$; per questo valore di $k$, calcolare 
                la firma $A$ del messaggio $P=15$.
            \item Verificare se anche la firma $A^{\prime}=(r^{\prime}, s^{\prime})=(12,69)$ è valida per 
                lo stesso messaggio $P=15$; se è valida, calcolare il valore di $k$ per cui è stata ricavata da Bob.
        \end{enumerate}


    \part{Soluzioni}
    %! TEX root = main.tex

\chapter{Aritmetica modulare}

    \bigskip
    \section{Potenze}
        \ASolution{11/07/2013}

        \ASolution{04/07/2014}

    \bigskip
    \section{Radici}
        \ASolution{27/02/2013}

        \ASolution{18/07/2014}

        \ASolution{31/08/2019}

    \bigskip
    \section{Inverso}
        \ASolution{08/02/2013}

        \ASolution{11/07/2013}

        \ASolution{04/07/2014}

        \ASolution{17/09/2014}

        \ASolution{03/07/2015}

    \bigskip
    \section{Residui quadratici}
        \ASolution{08/02/2013}

        \ASolution{27/02/2013}

        \ASolution{14/02/2014}

        \ASolution{17/09/2014}

        \ASolution{02/03/2015}

        \ASolution{03/07/2015}

        \ASolution{09/02/2017}

    \bigskip
    \section{Logaritmo discreto}
        \ASolution{08/02/2013}

        \ASolution{03/03/2014}

        \ASolution{17/07/2015}

        \ASolution{23/09/2016}

        \ASolution{06/07/2018}

    \bigskip
    \section{Fattorizzazione}
        \ASolution{04/07/2016}

        \ASolution{21/07/2016}

        \ASolution{09/02/2017}

        \ASolution{23/09/2016}

        \ASolution{30/08/2018}

        \ASolution{11/02/2019}

        \ASolution{26/07/2019}

        \ASolution{31/08/2019}


\end{document}
